\documentclass[11pt]{article}
\usepackage[T1]{fontenc}
\usepackage[utf8]{inputenc}
\usepackage{lmodern,microtype}
\usepackage[margin=1.08in]{geometry}
\usepackage{amsmath,amssymb,amsthm,mathtools}
\usepackage{xcolor,hyperref}
\hypersetup{colorlinks=true,linkcolor=blue!55!black,citecolor=blue!55!black,urlcolor=blue!60!black}
\setlength{\parindent}{0pt}
\setlength{\parskip}{0.48em}
\setlength{\emergencystretch}{2em}
\newtheorem{theorem}{Theorem}
\newtheorem{lemma}[theorem]{Lemma}
\newtheorem{proposition}[theorem]{Proposition}
\theoremstyle{remark}
\newtheorem{remark}[theorem]{Remark}

\title{Eventual Finiteness of Critical Steklov Lengths on Hypersurfaces of Revolution}
\author{[Author Name]}
\date{}

\begin{document}
\maketitle

\begin{abstract}
\textbf{Relation to the problem list.}
This manuscript addresses Question~38 (L38), ``Only finitely many infinite
critical Steklov lengths.''  It gives a full proof of the conjecture in its
sharp-bound-profile formulation: for each fixed dimension, every sufficiently
high index has finite critical length.  This is a statement about where the
sharp profile attains its supremum; it does not assert that a smooth maximizing
metric realizes the profile.

For hypersurfaces of revolution with two unit spherical boundary components,
M\'{e}tras and Tschanz associated to each Steklov index a sharp-bound profile
depending on the meridian length and conjectured that every sufficiently high
index has a finite critical length.  We prove the conjecture in every
dimension.  At a sequence of diagnosis indices, we choose the coupled scale
$1+L/2=e^{s/i}$ and count simultaneously the increasing Neumann and
decreasing Dirichlet branches.  The weighted count reduces to the scalar
inequality $a_s^{n-1}+b_s^{n-1}<2$, where
$a_s\tanh(sa_s)=1$ and $b_s\coth(sb_s)=1$.  A second-order large-$s$
expansion proves this inequality in every fixed dimension.  A propagation
lemma then covers all intervening indices.
\end{abstract}

\section{Introduction}

Let
\[
 M=[0,L]\times\mathbb S^{n-1},
 \qquad
 g=dr^2+h(r)^2g_{\mathbb S^{n-1}},
\]
where $L>0$, $h(0)=h(L)=1$, $h>0$, and $|h'|\le1$.  The two boundary
components are unit round spheres.  M\'{e}tras and Tschanz derived, for
each positive Steklov index $k$, a continuous sharp-bound profile
$B_n^k(L)$ and defined
\[
 B_n^k=\sup_{L>0}B_n^k(L).
\]
The index $k$ has \emph{finite critical length} if this supremum is attained
by the profile at some finite $L$, and it has \emph{infinite critical length}
if
\[
 B_n^k=\lim_{L\to\infty}B_n^k(L).
\]
The profile is a sharp envelope; no assertion that a smooth maximizing
metric realizes it is needed here.  Conjecture~8 of \cite{MetrasTschanz}
asks whether, for each dimension, every sufficiently high index has finite
critical length.

\begin{theorem}\label{thm:main}
For every integer $n\ge2$, there exists $K(n)$ such that every Steklov
index $k\ge K(n)$ has finite critical length.
\end{theorem}

The case $n=2$ follows from Fan, Tam, and Yu \cite{FanTamYu}.  We prove the
result for every fixed $n\ge3$ by a full weighted-branch count.

\section{Mixed branches and weighted order statistics}

Put $R=1+L/2$ and consider the Euclidean annulus
$A_R=B_R\setminus\overline B_1\subset\mathbb R^n$.  For spherical-harmonic
degree $j\ge0$, the mixed Steklov--Dirichlet and Steklov--Neumann eigenvalues
at the inner sphere are
\begin{align}
 D_j(R)&=
 \frac{(j+n-2)R^{2j+n-2}+j}{R^{2j+n-2}-1},\label{eq:D}\\
 N_j(R)&=
 j\frac{(j+n-2)(R^{2j+n-2}-1)}
 {jR^{2j+n-2}+j+n-2}.\label{eq:N}
\end{align}
The branch $D_j$ decreases in $R$, the branch $N_j$ increases in $R$, and
both have spherical multiplicity
\begin{equation}\label{eq:multiplicity}
 m_j=\frac{(2j+n-2)(j+n-3)!}{j!(n-2)!}.
\end{equation}
Their cumulative multiplicity is
\begin{equation}\label{eq:S}
 S_J=\sum_{j=0}^Jm_j
 =\binom{J+n-1}{n-1}+\binom{J+n-2}{n-1}.
\end{equation}

For fixed $n,k,L$, let $\ell_0=\max\{\ell:S_\ell\le k\}$.  The extension
process of \cite{MetrasTschanz} forms the weighted multiset
\begin{equation}\label{eq:multiset}
 \{D_0,\ldots,D_{\ell_0},N_1,\ldots,N_{\ell_0+1}\},
\end{equation}
where every occurrence of degree $j$ has weight $m_j$.  The number
$B_n^k(L)$ is its $k$th weighted order statistic.  Since
\eqref{eq:multiset} is finite and its branches are continuous,
$B_n^k(L)$ is continuous in $L$.

Set
\[
 q=n-1\ge2,
 \qquad
 k_i=2S_{i-1}.
\]
These are the diagnosis indices used in \cite{MetrasTschanz}.

\begin{lemma}\label{lem:infinity}
For all sufficiently large $L$,
\[
 B_n^{k_i}(L)=N_i(1+L/2),
\]
and therefore
\begin{equation}\label{eq:Ti}
 \lim_{L\to\infty}B_n^{k_i}(L)=T_i:=i+q-1.
\end{equation}
\end{lemma}

\begin{proof}
For each fixed $j$, formulas \eqref{eq:D}--\eqref{eq:N} give
\[
 D_j(R),N_j(R)\longrightarrow j+q-1
 \qquad(R\to\infty).
\]
For large $R$, the branches $D_0,\ldots,D_{i-1}$ contribute
$S_{i-1}$ weighted occurrences below $N_i$, while
$N_1,\ldots,N_{i-1}$ contribute $S_{i-1}-1$; the subtraction removes the
zero Neumann mode.  Thus exactly $k_i-1$ occurrences precede $N_i$, proving
the first assertion.  The limit of $N_i$ gives \eqref{eq:Ti}.
\end{proof}

It is enough to find one finite $L$ for which
$B_n^{k_i}(L)>T_i$.

\section{The coupled branch limit}

Fix $s>1$ and set
\begin{equation}\label{eq:scale}
 R_i=e^{s/i},
 \qquad
 L_i=2(e^{s/i}-1).
\end{equation}

\begin{lemma}\label{lem:limits}
If $j_i/i\to x\ge0$, then
\begin{align}
 \frac{N_{j_i}(R_i)}{i}&\longrightarrow
 f_s(x):=x\tanh(sx),\label{eq:flimit}\\
 \frac{D_{j_i}(R_i)}{i}&\longrightarrow
 g_s(x):=x\coth(sx),\label{eq:glimit}
\end{align}
where $g_s(0)=1/s$.  Both $f_s$ and $g_s$ are strictly increasing on
$[0,\infty)$.
\end{lemma}

\begin{proof}
Substitute $R_i=e^{s/i}$ and $j_i/i\to x$ into
\eqref{eq:D}--\eqref{eq:N}; then
$R_i^{2j_i+n-2}\to e^{2sx}$, giving
\eqref{eq:flimit}--\eqref{eq:glimit}.  At $x=0$, the Dirichlet formula
gives $D_0(R_i)/i\to1/s$.  Finally,
\[
 f_s'(x)=\tanh(sx)+sx\operatorname{sech}^2(sx)>0,
\]
while, with $z=sx$,
\[
 g_s'(x)=\coth z-z\operatorname{csch}^2z
 =\frac{\sinh z\cosh z-z}{\sinh^2z}>0.
\]
\end{proof}

Let $a_s>1$ and $b_s\in(0,1)$ be the unique roots
\begin{equation}\label{eq:roots}
 a_s\tanh(sa_s)=1,
 \qquad
 b_s\coth(sb_s)=1.
\end{equation}
The second equation is equivalent to $b_s=\tanh(sb_s)$.  Define the
discrete cutoffs
\begin{align}
 A_i&=\max\{j\ge1:N_j(R_i)\le T_i\},\label{eq:Ai}\\
 C_i&=\max\{j\ge0:D_j(R_i)\le T_i\}.\label{eq:Ci}
\end{align}

\begin{lemma}\label{lem:cutoffs}
For fixed $s>1$,
\[
 \frac{A_i}{i}\longrightarrow a_s,
 \qquad
 \frac{C_i}{i}\longrightarrow b_s.
\]
\end{lemma}

\begin{proof}
The mixed eigenvalues are strictly increasing in the spherical degree.
Since $T_i/i\to1$, Lemma~\ref{lem:limits} and the strict monotonicity of
$f_s,g_s$ give the conclusion by a two-sided root sandwich.  For example,
for every $\varepsilon>0$, uniform convergence on the two fixed compact
intervals around $a_s\pm\varepsilon$ places $A_i/i$ between those two
numbers for all large $i$; the argument for $C_i$ is identical.
\end{proof}

\section{Weighted rank and the scalar inequality}

Count all mixed branch occurrences at $R_i$ whose value is at most $T_i$.
By \eqref{eq:Ai}--\eqref{eq:Ci}, this weighted rank is
\begin{equation}\label{eq:rank}
 \mathcal R_i=S_{A_i}-1+S_{C_i}.
\end{equation}
The $-1$ again removes $N_0=0$.  From \eqref{eq:S},
\begin{equation}\label{eq:Sasymptotic}
 S_J=\frac{2}{q!}J^q+O_q(J^{q-1}).
\end{equation}
Consequently,
\begin{equation}\label{eq:ranklimits}
 \frac{\mathcal R_i}{i^q}
 \longrightarrow\frac{2}{q!}(a_s^q+b_s^q),
 \qquad
 \frac{k_i}{i^q}\longrightarrow\frac4{q!}.
\end{equation}

\begin{lemma}\label{lem:scalar}
For every fixed integer $q\ge2$, there exists $s>1$ such that
\begin{equation}\label{eq:scalar}
 a_s^q+b_s^q<2.
\end{equation}
\end{lemma}

\begin{proof}
Put
\[
 E=e^{-2s},
 \qquad
 a_s=1+\delta_s,
 \qquad
 b_s=1-\varepsilon_s.
\]
Equations \eqref{eq:roots} are equivalent to
\begin{equation}\label{eq:exactroots}
 \frac{\delta_s}{2+\delta_s}=Ee^{-2s\delta_s},
 \qquad
 \frac{\varepsilon_s}{2-\varepsilon_s}=Ee^{2s\varepsilon_s}.
\end{equation}
The relevant roots satisfy $\delta_s,\varepsilon_s=O(E)$.  Indeed,
evaluating the defining functions at $1+4E$ and $1-4E$, and using the
elementary bounds for $\tanh$, gives
$0<\delta_s,\varepsilon_s<4E$ for all large $s$.

Taylor expansion of \eqref{eq:exactroots}, with $sE\to0$, yields
\begin{align}
 \delta_s&=2E+(2-8s)E^2+O(s^2E^3),\label{eq:deltaexp}\\
 \varepsilon_s&=2E+(8s-2)E^2+O(s^2E^3).\label{eq:epsilonexp}
\end{align}
For fixed $q$, the binomial expansion gives
\begin{align}
 a_s^q+b_s^q-2
 &=q(\delta_s-\varepsilon_s)
 +\frac{q(q-1)}2(\delta_s^2+\varepsilon_s^2)+O_q(E^3)\notag\\
 &=4q(q-4s)E^2+O_q(s^2E^3).\label{eq:sign}
\end{align}
The leading term is asymptotic to $-16qsE^2$, whereas the remainder divided
by $sE^2$ tends to zero.  Hence \eqref{eq:scalar} holds for all sufficiently
large $s$.
\end{proof}

\begin{proposition}\label{prop:diagnosis}
For every fixed $n\ge3$, all sufficiently large diagnosis indices $k_i$
have finite critical length.
\end{proposition}

\begin{proof}
Fix $q=n-1$ and choose $s$ as in Lemma~\ref{lem:scalar}.  Equations
\eqref{eq:ranklimits} then give $\mathcal R_i<k_i$ for all sufficiently
large $i$.  The finite multiset \eqref{eq:multiset} is a sub-multiset of the
full family counted in \eqref{eq:rank}.  Hence fewer than $k_i$ of its
weighted occurrences can be at most $T_i$, and therefore
\begin{equation}\label{eq:overshoot}
 B_n^{k_i}(L_i)>T_i.
\end{equation}
By Lemma~\ref{lem:infinity}, the right-hand side is the infinite-length
limit of the profile.

For fixed $i$, the function $B_n^{k_i}(L)$ is continuous, tends to zero as
$L\downarrow0$, and has the finite limit $T_i$ as $L\to\infty$.
The strict overshoot \eqref{eq:overshoot} confines its global supremum to a
compact subinterval of $(0,\infty)$.  Thus the supremum is attained at a
finite length.
\end{proof}

\section{Propagation to every high index}

We use the following propagation result of M\'{e}tras and Tschanz
\cite[Lemma~17]{MetrasTschanz}.

\begin{lemma}[Block propagation]\label{lem:propagation}
If $k_i=2S_{i-1}$ has finite critical length, then every index
\[
 k_i\le k\le k_i+2m_i-1=k_{i+1}-1
\]
has finite critical length.
\end{lemma}

\begin{proof}[Proof of Theorem~\ref{thm:main}]
For $n\ge3$, Proposition~\ref{prop:diagnosis} gives an $i_0$ such that
every $k_i$ with $i\ge i_0$ has finite critical length.  Applying
Lemma~\ref{lem:propagation} to each such $i$ covers the consecutive blocks
$[k_i,k_{i+1}-1]$.  Their union is every index $k\ge k_{i_0}$.  Thus one
may take $K(n)=k_{i_0}$.  The two-dimensional case is proved in
\cite{FanTamYu}, completing the theorem.
\end{proof}

\begin{remark}
The choice of $s$ depends on $q=n-1$.  The gap in
\eqref{eq:scalar} becomes small in high dimension, so the argument is
dimension-by-dimension and does not provide a useful uniform bound for
$K(n)$.
\end{remark}

\begin{thebibliography}{9}

\bibitem{MetrasTschanz}
A.~M\'{e}tras and L.~Tschanz,
\emph{Critical lengths of Steklov eigenvalues of hypersurfaces of revolution
in Euclidean space},
Experimental Math. \textbf{34} (2025), 728--744;
\href{https://doi.org/10.1080/10586458.2024.2410967}{DOI}.

\bibitem{Tschanz}
L.~Tschanz,
\emph{Sharp upper bounds for Steklov eigenvalues of a hypersurface of
revolution with two boundary components in Euclidean space},
Ann. Math. Qu\'{e}bec \textbf{48} (2024);
\href{https://doi.org/10.1007/s40316-024-00225-8}{DOI}.

\bibitem{FanTamYu}
X.-Q.~Fan, L.-F.~Tam, and C.~Yu,
\emph{Extremal problems for Steklov eigenvalues on annuli},
Calc. Var. Partial Differential Equations \textbf{54} (2015), 1043--1059;
\href{https://doi.org/10.1007/s00526-014-0816-8}{DOI}.

\end{thebibliography}
\end{document}
